\subsection{原函数的主部}

设 \(f\) 是一个规则实函数， \(\neq 0\) 且在一个区间 \([a, +\infty[\) 上具有常号；下面的命题使我们在某些情形能够得到原函数 \(\int_{x}^{+\infty} f(t) dt\) 的主部（若 \(\int_{a}^{+\infty} f(t) dt\) 收敛），以及原函数 \(\int_{a}^{x} f(t) dt\) 的主部（若 \(\int_{a}^{+\infty} f(t) dt\) 为无穷）：

命题 8. 置 \(\mathrm{F}(x) = \int_{x}^{+\infty} f(t) dt\) （若 \(\int_{a}^{+\infty} f(t) dt\) 收敛），而置 \(\mathrm{F}(x) = \int_{a}^{x} f(t) dt\) （若 \(\int_{a}^{+\infty} f(t) dt\) 为无穷）。设 \(\log |f|\) 与 \(\log x\) 是 1 阶可比较的。

\(1^{\circ}\) 若 \(f\) 是有限阶 \(\mu \neq -1\) 的（相对于 \(x\) ），则有

\[
\mathrm{F} (x) \sim \frac {1}{| \mu + 1 |} x f (x).\tag{1}
\]

\(2^{\circ}\) 若 \(f\) 相对于 \(x\) 是无穷阶的，且 \(f / f'\) 与 \(x\) 是 1 阶可比较的，则有

\[
\mathrm{F} (x) \sim \frac {(f (x)) ^ {2}}{| f ^ {\prime} (x) |}.\tag{2}
\]

应当注意，假设蕴含 \(f(x)\) 相对于 \(x\) 有确定的阶 (V, p.~219)。

\(1^{\circ}\) 若 \(f\) 是 \(\mu \neq 0\) 阶的（相对于 \(x\) ），则有 \(\log |f| \sim \mu \log x\) ，于是，由于 \(\log |f|\) 与 \(\log x\) 是 1 阶可比较的，由 V, p.~232 的命题 7 有 \(f'/f \sim\) \(\mu/x\) ，由此 \(xf' \sim \mu f\) 。若 \(\mu > -1\) ，则有 \(f(x) \gg x^{\mu-\varepsilon}\) （对每个 \(\varepsilon > 0\) ），从而 (V, p.~229, prop. 2) 积分 \(\int_{a}^{+\infty} f(t) dt\) 为无穷。可以写 \(F(x) = \int_{a}^{x} f(t) dt = xf(x) - af(a) - \int_{a}^{x} tf'(t) dt\) ，由此再次得到

\[
\int_ {a} ^ {x} \left(f (t) + t f ^ {\prime} (t)\right) d t = x f (x) - a f (a);
\]

由于 \(\mu \neq -1\) ，有 \(f(x) + xf'(x) \sim (\mu + 1)f(x)\) ，故 (V, p.~230, prop. 6)

\[
\int_ {a} ^ {x} \left(f (t) + t f ^ {\prime} (t)\right) d t \sim (\mu + 1) \mathrm{F} (x),
\]

这就在这种情形证明了 (1)。若 \(\mu = 0\) ，同样有 \(xf'(x) \ll f(x)\) ，它再次给出 \(f(x) + xf'(x) \sim f(x)\) 。当 \(\mu < -1\) 时论证类似（在 \(\int_{a}^{+\infty} f(t) dt\) 收敛的情形）。

\(2^{\circ}\) 若 \(f\) 是 \(+\infty\) 阶的（相对于 \(x\) ），则有 \(\log |f| \gg \log x\) ，故 (V, p.~232, prop. 7) \(f'/f \gg 1/x\) ，或者再令 \(g(x) = f(x)/f'(x)\) ，有 \(g(x) \ll x\) ；此外，由于 \(f(x) \gg x^{\alpha}\) （对每个 \(\alpha > 0\) ），积分 \(\int_{a}^{+\infty} f(t) dt\) 为无穷。可以写

\[
\mathrm{F} (x) = \int_ {a} ^ {x} f (t) d t = \int_ {a} ^ {x} g (t) f ^ {\prime} (t) d t = g (x) f (x) - g (a) f (a) - \int_ {a} ^ {x} f (t) g ^ {\prime} (t) d t;
\]

由于 \(g\) 与 \(x\) 是 1 阶可比较的，从关系 \(g(x) \ll x\) 可推出 (V, p.~232, prop. 7) \(g'(x) \ll 1\) ，从而 \(fg' \ll f\) ，于是 (V, p.~230, prop.6)

\[
\int_ {a} ^ {x} f (t) g ^ {\prime} (t) d t \ll \mathrm{F} (x),
\]

这就建立了关系式 (2)。当 \(f\) 是 \(-\infty\) 阶的（相对于 \(x\) ）时证明类似（在 \(\int_{a}^{+\infty} f(t) dt\) 收敛的情形）。

设 E 是一个比较尺度（对趋于 \(+\infty\) 的实 x 而言），由在 \(+\infty\) 的一个邻域上具有常号的非零实函数组成，使得 \(x \in E\) ，且 E 中两个函数的乘积与商仍属于 E (V, p.~221 and p.~224)。若一个在 \(+\infty\) 的一个邻域上具有常号的规则函数 f 相对于 E 有主部 cg，则 \(\int_{x}^{+\infty} f(t) dt\) （相应地， \(\int_{a}^{x} f(t) dt\) ，视情形而定）将等价于 \(c \int_{x}^{+\infty} g(t) dt\) （相应地， \(c \int_{a}^{x} g(t) dt\) ）；若函数 g 满足 V, p.~233 的命题 8 的假设，且（当 V, p.~233 的公式 (2) 适用时）知道 \(g'\) 相对于 E 的一个主部，则由此将有 \(\int_{x}^{+\infty} f(t) dt\) （相应地， \(\int_{a}^{x} f(t) dt\) ）相对于 E 的一个主部。

例. 1) 函数 \(1 / \log x\) 相对于 \(x\) 是 0 阶的且满足 V, p.~233 的命题 8 的条件；于是

\[
\int_ {a} ^ {x} \frac {d t}{\log t} \sim \frac {x}{\log x}.
\]

\begin{enumerate}
\def\labelenumi{\arabic{enumi})}
\setcounter{enumi}{1}
\tightlist
\item
  函数 \(e^{x^2}\) 是 \(+\infty\) 阶的（相对于 \(x\) ）且满足命题 8 的条件，故
\end{enumerate}

\[
\int_ {a} ^ {x} e ^ {t ^ {2}} d t \sim \frac {1}{2 x} e ^ {x ^ {2}}.
\]

在附录 (V, p.~252) 中我们将定义一类函数，对它们命题 7 与命题 8 总是适用的。

注. 命题 8 不直接应用于一个函数 \(f\) （ \(-1\) 阶的，相对于 \(x\) ）。但这时可以写 \(f(x) = f_1(x) / x\) ，其中 \(f_1\) 相对于 \(x\) 是 0 阶的。例如设 \(\int_{a}^{+\infty} f(t) dt\) 为无穷；则

\[
\mathrm{F} (x) = \int_ {a} ^ {x} f (t) d t = \int_ {a} ^ {x} \frac {1}{t} f _ {1} (t) d t = \int_ {\log a} ^ {\log x} f _ {1} \left(e ^ {u}\right) d u.
\]

若函数 \(f_{1}(e^{y})\) 满足命题 8 的假设且有一个 \(\neq -1\) 的阶（相对于 \(y\) ）（也就是说，若 \(f_{1}(x)\) 有一个 \(\neq -1\) 的阶（相对于 \(\log x\) ）），则公式 (1) 与 (2) 仍能使我们得到 \(F(x)\) 的一个主部。例如，设 \(f(x) = \frac{\exp(\sqrt{\log x})}{x\log x}\) ；由于 \(\exp(\sqrt{\log x})\) 相对于 \(x\) 是 0 阶的， \(f\) 是 \(-1\) 阶的；这里有 \(f_{1}(e^{y}) = e^{\sqrt{y}} / y\) ，且这个函数是 \(+\infty\) 阶的（相对于 \(y\) ）；命题 8 适用并给出 \(\int_{\alpha}^{y} e^{\sqrt{u}} / u du \sim 2e^{\sqrt{y}} / \sqrt{y}\) ；回到变量 \(x\) ，我们有 \(\int_{a}^{x} \frac{\exp(\sqrt{\log t})}{t\log t} dt \sim \frac{2\exp(\sqrt{\log x})}{\sqrt{\log x}}\) 。

